\section{Discussion and Limitations}
\subsection{What the pose experiments establish}
The corrected crowd study gives a modest mean accuracy advantage to interactions and reliability relative to the baseline, with no combined gain. This outcome does not establish a universal benefit or a failure of interaction modeling in general. The nearest-person features may be useful in some recordings and noisy in others. Added confidence channels can provide information while increasing parameter count and exposing estimator behavior. The current experiment does not isolate which of these possibilities drives the combined result.

Centering changes matter because derived velocities depend on the coordinate origin. The matched CPU comparison isolates that construction choice more carefully than comparing historical GPU fitting with a corrected CPU run. Its mixed variant results are a reason to record preprocessing precisely. A stable origin does not repair all missing joints, wrong associations, or camera motion, and complete-segment interpolation prevents a direct causal streaming interpretation.

The detection controls also require precise language. Competitive quality-only performance demonstrates predictive aggregate cues. Detection-only TCN performance shows that retained per-frame quality and presence information remains useful after zeroing selected kinematics. Neither proves that every full-model decision depends on a shortcut or identifies a specific scene bias. The unchanged explicit-coverage means cannot eliminate this concern because indirect cues remain. The ordering control retains precomputed velocity and therefore cannot support the statement that motion is irrelevant.

\subsection{Interpreting the video comparison}
The frozen crowd video reference provides a stronger result under the same clip assignments, but it has substantial action pretraining and different model capacity. Pretraining-source overlap is unaudited. Its result motivates a practical reference and further testing, rather than a conclusion that pose features are inherently inferior across datasets, architectures, or tasks.

The supervised RWF reference changes the downstream training domain. The later fine-tuning comparison is within that domain and follows an explicit validation rule. Its higher AUC and Fight recall accompany lower precision and slightly lower accuracy on RWF. The hockey result improves both accuracy and recall but still has 132 missed fights and 57 false alarms. Thus improved external accuracy does not mean every relevant operating requirement is met.

A lower selected threshold can account for part of the recall/false-alarm shift. Separating the effect of adapted weights from operating-point selection would require additional matched-threshold or fixed-cost analyses declared independently of the completed evaluations. We do not retroactively tune the current threshold to obtain a more favorable comparison. A deployment setting would also need a justified balance between missed incidents and unnecessary review; macro-F1 on development clips is not itself that application policy.

\subsection{Data exposure and uncertainty}
The crowd and RWF footage had been inspected before the final studies. Source grouping and later sealing improve internal separation, but they do not restore an untouched test. Known source merges and duplicate exclusions address identified overlap, while unknown reused footage and Kinetics overlap remain possible. Hockey's byte comparison rules out byte-identical files against the declared inventory but does not establish independent games or complete near-duplicate separation.

Labels remain supplied and independently unreviewed. A high-confidence false positive could reflect classifier failure, label ambiguity, or both; a selected frame cannot distinguish those explanations. A future independent review should assess full clips under a documented rubric before inspecting model scores. That process should preserve the current results and record any revised-label experiment separately rather than silently changing the completed evaluation.

Fold/seed variation is descriptive. Repeated clips, shared training sets, and known source dependence invalidate an interpretation of all repeated decisions as independent observations. The manuscript therefore reports no significance claim from small mean differences. The fine-tuning result also comes from one bounded seed-42 run. It does not establish robustness over fine-tuning seeds, hardware, or hyperparameter budgets.

\subsection{Scope, privacy, and operational behavior}
All primary outputs concern complete short clips. Uniform sparse sampling can miss short events, average pooling can lose temporal position, and masked person pooling can lose individual attribution. A closed-window interface can process longer videos offline, but window accounting alone does not validate outbreak detection, continuous-stream false-alarm rates, or latency under live load. The pose benchmark covers ten fixed clips and one laptop GPU, with export excluded.

Pose representation can omit some raw appearance from classifier input, yet the pipeline still processes videos and can expose contextual or identifying motion. It implements no formal privacy guarantee, federated system, or security audit. Neither tracks nor clip scores establish responsibility, intent, or individual guilt. Any operational proposal would require additional reviewed data, domain-specific validation, error-cost analysis, and handling of uncertain or ambiguous clips.

The next research step is a fresh, independently reviewed evaluation corpus with more complete recording-source information. Predeclared analyses could examine fixed operating points, stability over seeds, actual occlusion and association failures, and temporal-event coverage. Those are future experiments. The current archive supplies reproducible evidence of implemented behavior and its limitations, including negative findings, rather than substituting computational checks for independent evaluation.
